\subsection*{关系式已经足够描述曲线}
圆由 $x^2+y^2=1$ 描述，不必先把 $y$ 解成一个全局函数。沿圆移动时，$x$ 和 $y$ 不能自由独立变化，它们必须共同维持这个约束。若 $x$ 增加一点，另一个坐标要怎样配合，才能仍在曲线上？

对关系作一阶变化，得到 $2x\,dx+2y\,dy=0$。这里的 $dy$ 不是任意小量，而是曲线约束要求的响应。若 $y\ne0$，解出 $dy/dx=-x/y$，就得到局部斜率，无须展开平方根。

\subsection*{斜率不存在，不代表切线不存在}
在 $(3/5,4/5)$ 处，斜率为 $-3/4$。向右走一点必须略向下，才能补偿横坐标平方的增加。切线给出这种补偿的一阶近似；有限移动后的点通常不在圆上，因为曲线还存在二阶弯曲。

来到 $(1,0)$，公式不能除以 $y$。这并非圆突然不光滑，而是此处不能用单值的 $y(x)$ 斜率描述垂直切线。改用 $x$ 作为 $y$ 的局部函数，就能恢复合适的坐标语言。一个几何对象是否规则，与某种表示是否方便，是两个问题。

一般关系 $F(x,y)=0$ 中，沿曲线总变化为零，得到 $F_x+F_y y'=0$。使用它时仍须核对相关可微性及局部函数存在条件。Lean 对照只验证微分约束后的代数解法；不要把这一点误认为已经证明任意隐式关系都能解成光滑函数。
